\documentclass[conference]{IEEEtran}
\IEEEoverridecommandlockouts
\usepackage{cite}
\usepackage{amsmath,amssymb,amsfonts}
\usepackage{graphicx}
\usepackage{textcomp}
\usepackage{url}
\usepackage{caption}
\usepackage{amsmath,amssymb}
\usepackage{algorithm}
\usepackage{algpseudocode}

\begin{document}

\title{Smart Grid Sensor Network for Industrial Energy\\
Management Using IoT and Machine Learning}

\author{\IEEEauthorblockN{Saeed Sarwar Anas}
\IEEEauthorblockA{\textit{School of Computing, Engineering and Intelligent Systems}\\
\textit{Ulster University}\\
Student Number: 20068400\\
Supervisor: Dr Anwar Haq}}

\maketitle
\begin{abstract}
Industrial facilities account for a substantial share of global electricity
demand, yet most small and medium-sized enterprises (SMEs) operate without any
real-time visibility of their energy use because commercial monitoring
infrastructure is prohibitively expensive. This paper presents the design and
evaluation of a low-cost smart grid sensor network for industrial energy
management that integrates Internet of Things (IoT) sensing with machine
learning. The system couples an inexpensive embedded sensing platform with a
reproducible data pipeline, a Long Short-Term Memory (LSTM) model for
short-term demand forecasting, and an Isolation Forest model for unsupervised
anomaly detection, all surfaced through a real-time dashboard. Using the UCI
Individual Household Electric Power Consumption dataset, the LSTM forecaster
achieved a mean absolute error (MAE) of 0.3707~kW, outperforming a naive
baseline by 8.8\% and a 24-hour moving-average baseline by 34.9\% on unseen
test data. The contribution is an affordable, open-source, end-to-end pipeline
that lowers the barrier to intelligent energy management for resource-
constrained organisations.
\end{abstract}

\begin{IEEEkeywords}
smart grid, Internet of Things, machine learning, LSTM, energy demand
forecasting, anomaly detection, isolation forest, edge computing, Industry~4.0
\end{IEEEkeywords}

\section{Introduction}
The electricity grid is among the most important engineering systems of the
modern era, but for most of the twentieth century it operated as an essentially
``one-way'' network. Electricity flowed from a small number of large, central
generators to passive consumers, and the grid carried almost no information back
about how, when, or where energy was being used. This architecture was robust
and well suited to the demand patterns of its time, yet it left operators
fundamentally blind to consumption at the point of use and unable to respond to
changing conditions in real time.

Several converging pressures over the past two decades have exposed the limits
of this traditional model. Global electricity demand has continued to rise, with
industry responsible for a large fraction of total consumption
\cite{international2022international}. At the same time, decarbonisation targets
and the rapid uptake of intermittent renewable generation have demanded a far
more flexible and observable grid. These pressures gave rise to the concept of
the \emph{smart grid}: an electricity network augmented with digital sensing,
two-way communication, and data-driven control, so that supply and demand can be
monitored and balanced dynamically rather than managed reactively. Surveys of
smart grid communication infrastructure document this shift from a passive
network towards a sensor-rich, communication-enabled system
\cite{fan2012smart,yan2012survey}.

Three broad paradigms have shaped how the sector has evolved since then. The
first is the move from centralised generation towards \emph{distributed} and
decentralised energy resources, which multiplies the number of points that must
be observed and coordinated \cite{silva2023iot}. The second is the rise of
\emph{Industry~4.0} and the Industrial Internet of Things (IIoT), in which
low-cost networked sensors make pervasive, fine-grained measurement economically
feasible \cite{munirathinam2020industry}. The third is the application of
\emph{big-data analytics and machine learning}, which has allowed energy systems
to move from descriptive, threshold-based monitoring towards predictive and
proactive management. Deep learning in particular has proven effective at
modelling complex, non-linear patterns in sequential data \cite{lecun2015deep},
and recurrent architectures such as the Long Short-Term Memory (LSTM) network
\cite{hochreiter1997long} are well matched to the daily, weekly, and seasonal
cycles that characterise energy consumption. In parallel, unsupervised
techniques such as the Isolation Forest \cite{liu2008isolation} have made it
possible to flag abnormal consumption without requiring labelled fault data,
which is rarely available in practice.

The combined impact of these approaches is significant. Sensor-driven monitoring
of buildings and homes has been shown to reduce energy consumption through
real-time feedback and automated control \cite{alaa2017review,hafeez2020efficient};
machine-learning forecasters, and LSTM models in particular, have been shown to
outperform classical statistical baselines on short-term load prediction
\cite{kong2017short,raza2015review}; and anomaly-detection methods have been used
to identify equipment faults and energy waste before they escalate
\cite{himeur2021artificial}. Together they point towards a future in which energy
is monitored intelligently and continuously at the point of use.

Despite this progress, a clear gap remains. The vast majority of deployed smart
grid and industrial monitoring systems rely on expensive proprietary
infrastructure---supervisory control and data acquisition (SCADA) systems,
industrial programmable logic controllers, and commercial smart meters---whose
cost runs to many thousands of pounds. This places intelligent energy management
firmly out of reach of small and medium-sized enterprises (SMEs), even though
SMEs collectively account for a large share of industrial energy use. Moreover,
much of the existing research addresses individual components in isolation---a
forecasting model here \cite{wen2024deep}, an anomaly detector there
\cite{pereira2018unsupervised}---rather than demonstrating a complete,
reproducible pipeline. As a result, an affordable and openly reproducible
end-to-end system that an under-resourced organisation could actually adopt has
remained largely unaddressed.

This paper addresses that gap. It presents the design, implementation, and
evaluation of a low-cost smart grid sensor network for industrial energy
management that integrates IoT sensing with machine learning in a single,
reproducible pipeline. In contrast to prior work that depends on costly
proprietary hardware or that treats forecasting and anomaly detection
separately, the proposed system combines inexpensive embedded sensing (targeting
hardware of the order of \pounds30), a documented and version-controlled data
pipeline, an LSTM-based demand forecaster, an Isolation~Forest anomaly detector,
and a real-time monitoring dashboard. The work is grounded in a Design Science
Research methodology \cite{hevner2004design,peffers2007design} and follows
established guidance for the conduct of the research project
\cite{zina2021essential}.

\noindent The principal contributions of this paper are:
\begin{itemize}
  \item an affordable, open-source, end-to-end architecture that brings together
        low-cost IoT sensing, wireless telemetry, time-series forecasting,
        unsupervised anomaly detection, and real-time visualisation;
  \item an empirical evaluation of an LSTM demand forecaster on a public
        benchmark dataset, including a comparison against naive and
        moving-average baselines; and
  \item a reproducible pipeline---with code, processed data, and trained models
        released publicly---that lowers the barrier to intelligent energy
        management for SMEs.
\end{itemize}

The remainder of this paper is organised as follows. Section~\ref{sec:problem}
defines the problem and its scale. Section~\ref{sec:aims} states the aim and
objectives. Section~\ref{sec:litreview} reviews related work, critically
compares existing approaches, and derives the set of parameters that informs the
proposed solution architecture. Section~\ref{sec:method} describes the
methodology and system design, Section~\ref{sec:results} reports and discusses
the results, and Section~\ref{sec:conclusion} concludes and outlines future
work.

\section{Problem Statement}\label{sec:problem}
 
Industrial activity is one of the largest consumers of energy worldwide,
accounting for approximately one third of total global energy use
\cite{international2022international}.
% TODO: if you want a precise percentage here, insert the exact verified figure
% from the IEA report you are citing, and adjust the wording accordingly.
A substantial part of this energy is consumed inefficiently, and in many
facilities that inefficiency is never detected because consumption is not
measured at the level of individual processes or machines. The economic and
environmental cost of this undetected waste is considerable, and it grows as
global demand continues to rise and as decarbonisation targets place increasing
pressure on energy efficiency.
 
This burden falls disproportionately on small and medium-sized enterprises
(SMEs). SMEs constitute the majority of industrial and commercial organisations,
yet most operate without any form of intelligent energy monitoring. The
principal barrier is cost: established solutions---supervisory control and data
acquisition (SCADA) systems, industrial programmable logic controllers, and
commercial smart-metering platforms---typically require capital investment
running to tens of thousands of pounds, alongside specialist installation and
maintenance \cite{vijayalakshmi2023smart,silva2023iot}. For an organisation
operating on tight margins, such systems are simply not viable. The consequence
is that these businesses operate effectively ``blind'': energy waste accumulates
without being seen, equipment faults develop without early warning, and there is
no basis on which to forecast demand or optimise consumption.
 
Where monitoring does exist, it is frequently reactive. Conventional
threshold-based alarms register a problem only after a measured value has already
crossed a fixed limit, and they are poorly suited to detecting the gradual drift
or subtle, context-dependent anomalies that often precede equipment failure or
signal inefficiency \cite{himeur2021artificial}. Recent advances in IoT sensing
\cite{alaa2017review,hafeez2020efficient} and in machine-learning-based
forecasting \cite{kong2017short,raza2015review} indicate that a more proactive
approach is achievable; however, these techniques have largely been demonstrated
as isolated components rather than combined into a single deployable system.
 
The core problem this project addresses is therefore the absence of an
affordable, integrated, and reproducible system that can sense industrial
environmental conditions in real time, forecast short-term energy demand, and
automatically detect anomalous consumption, all within a hardware and software
budget accessible to an SME. The problem is non-trivial at this level because a
credible solution must integrate several distinct disciplines---embedded sensing
and wireless communication, time-series deep learning for forecasting,
unsupervised learning for anomaly detection, and data engineering for a
reproducible pipeline---into one coherent system, and must then be evaluated
rigorously against established baselines rather than merely asserted to work.
 
\section{Aim and Objectives}\label{sec:aims}
 
\textit{Aim.} The aim of this project is to design, implement, and evaluate a
low-cost smart grid sensor network for industrial energy management that uses
the Internet of Things and machine learning to deliver real-time monitoring,
short-term demand forecasting, and anomaly detection, and thereby to demonstrate
that capabilities previously confined to expensive proprietary systems can be
achieved with affordable hardware and open-source software.
 
\noindent The objectives that operationalise this aim are as follows:
 
\begin{enumerate}
\item To critically review the literature on IoT-based energy monitoring,
machine-learning demand forecasting, and anomaly detection, in order to identify
the research gap and define a precise project scope.
 
\item To acquire and prepare benchmark energy data---principally the UCI
Individual Household Electric Power Consumption dataset
\cite{hebrail2012individual}---through a documented, version-controlled
preprocessing pipeline implemented in Python with \texttt{pandas}, applying
chronological train/validation/test partitioning to prevent data leakage.
 
\item To design the IoT sensing architecture around a low-cost embedded platform, using Bluetooth Low
Energy and Wi-Fi telemetry to stream environmental measurements to a time-series
datastore \cite{munirathinam2020industry,warden2019tinyml}.
 
\item To develop and train a Long Short-Term Memory (LSTM) network
\cite{hochreiter1997long} for short-term energy demand forecasting using
TensorFlow/Keras \cite{abadi2016tensorflow}, and to benchmark it against naive
and moving-average baselines \cite{makridakis2018statistical}.
 
\item To implement unsupervised anomaly detection using an Isolation Forest model
\cite{liu2008isolation} that flags abnormal consumption patterns without
requiring labelled fault data.

\item To build a real-time monitoring dashboard using Streamlit that visualises
live readings, forecasted demand, and anomaly alerts in a form usable by a
non-specialist operator.
 
\item To evaluate all models rigorously on held-out test data using appropriate
quantitative metrics---mean absolute error (MAE), root-mean-square error (RMSE),
and mean absolute percentage error (MAPE) for forecasting, and precision,
recall, and F1-score for anomaly detection.
 
 
\end{enumerate}

\section{Literature Review}\label{sec:litreview}
Building intelligent energy management on low-cost hardware draws on four
overlapping bodies of work: IoT-based energy monitoring, machine-learning demand
forecasting, anomaly detection, and low-cost embedded and edge computing. This
section reviews each in turn, drawing out comparative strengths and limitations,
and then synthesises a set of design parameters that informs the proposed
solution architecture.
 
\subsection{IoT-Based Energy Monitoring}
The integration of the Internet of Things into energy systems has been surveyed
extensively. Goudarzi \textit{et al.} \cite{goudarzi2022survey} and Kirmani
\textit{et al.} \cite{kirmani2022survey} both provide broad surveys of
IoT-enabled smart grids, cataloguing enabling technologies, architectures, and
open challenges; together they establish that sensing, communication, and
analytics are now treated as a single integrated stack rather than separate
concerns. At the application level, reviews of smart-home IoT
\cite{alaa2017review} and of smart-meter data analysis \cite{yildiz2017recent}
show that fine-grained measurement combined with feedback can meaningfully
influence consumption, while Hafeez \textit{et al.} \cite{hafeez2020efficient}
couple IoT sensing with price-based demand response to actively manage household
load.
 
Moving from the home to larger systems, several authors present complete energy
management platforms. Silva \textit{et al.} \cite{silva2023iot} target AC
microgrids with explicit grid and security constraints; Rehman \textit{et al.}
\cite{rehman2022secure} foreground the security of an edge-based management
protocol; and Vijayalakshmi \textit{et al.} \cite{vijayalakshmi2023smart}
demonstrate energy management on IoT-based embedded systems. The work closest to
a complete pipeline is that of Bin Mofidul \textit{et al.} \cite{bin2022real},
which combines real-time data acquisition, anomaly detection, and monitoring
within an Industrial IoT infrastructure spanning both edge and cloud AI.
 
Two limitations recur across this literature. First, attention is concentrated
at the two ends of the scale spectrum---individual residences on one side and
utility-scale microgrids on the other---leaving the industrial SME, which lacks
both the simplicity of a home and the resources of a utility, comparatively
underserved. Second, the more capable integrated systems, such as
\cite{bin2022real}, achieve their breadth through relatively heavy and costly
edge--cloud infrastructure, which reintroduces precisely the affordability
barrier that constrains SMEs. The literature therefore motivates an integrated
system but leaves open the question of doing so cheaply and reproducibly.
 
\subsection{Machine Learning for Energy Demand Forecasting}
Demand forecasting underpins proactive energy management, and the field has
shifted decisively from classical statistical models towards deep learning.
Foundational advances in deep learning \cite{lecun2015deep,goodfellow2016deep},
the convolutional architectures introduced by LeCun \textit{et al.}
\cite{lecun1998gradient}, and the Long Short-Term Memory (LSTM) network of
Hochreiter and Schmidhuber \cite{hochreiter1997long} provide the methodological
basis for modelling the non-linear, multi-scale temporal patterns that
characterise energy data. Reviews by Raza and Khosravi \cite{raza2015review} and
Aguiar-P\'erez and P\'erez-Ju\'arez \cite{aguiar2023insight} trace the rapid
adoption of these methods for load and demand forecasting.
 
Within this space, the LSTM has become the dominant choice for short-term load
forecasting. Kong \textit{et al.} \cite{kong2017short} applied an LSTM
recurrent network to short-term residential load and demonstrated its ability to
capture household consumption dynamics; Inala \textit{et al.}
\cite{inala2025enhancing} apply the same family of models to household data in a
smart-grid scenario. Comparative studies refine this picture: Mahjoub
\textit{et al.} \cite{mahjoub2022predicting} benchmark LSTM against multi-layer
and drop-GRU variants, highlighting an accuracy--efficiency trade-off between
architectures, while Wen \textit{et al.} \cite{wen2024deep} pursue higher
accuracy through a hybrid deep model at the cost of additional complexity.
 
A crucial and often-overlooked caution comes from Makridakis \textit{et al.}
\cite{makridakis2018statistical}, who found that machine-learning forecasters do
not automatically outperform simpler statistical methods and argued for rigorous
evaluation against strong baselines. This finding directly shapes the present
work: rather than assuming the superiority of the LSTM, the forecaster is
benchmarked against naive and moving-average baselines on held-out data. The
broader gap in this literature is that most studies optimise predictive accuracy
on a single dataset and report little on the model's deployability---its size,
inference cost, or suitability for constrained hardware---which is precisely the
constraint that matters for low-cost deployment.
 
\subsection{Anomaly Detection in Energy Systems}
Detecting abnormal consumption is the second analytical pillar. Himeur
\textit{et al.} \cite{himeur2021artificial} review artificial-intelligence-based
anomaly detection for energy consumption in buildings, providing a taxonomy of
techniques and noting the recurring practical difficulty that labelled anomalies
are seldom available. This label scarcity is the central tension in the field
and divides the literature into two camps.
 
The first camp applies supervised or ensemble classifiers. Chen \textit{et al.}
\cite{chen2024effective} combine XGBoost and LSTM for anomaly detection in the
Industrial IoT; Aly and Behiry \cite{aly2025enhancing} compare logistic
boosting, random forest, and support vector machines for IoT-driven factories;
and Sahu and Mukherjee \cite{sahu2020machine} apply machine learning to IoT
network anomaly detection. These approaches can be highly accurate but depend on
labelled examples of the very faults they aim to catch, which limits their
applicability in real deployments where such labels do not exist.
 
The second camp uses unsupervised methods that require no labels. The Isolation
Forest of Liu \textit{et al.} \cite{liu2008isolation} isolates anomalies
efficiently and scales well; Jesmeen \textit{et al.} \cite{jesmeen2021unsupervised}
apply clustering to energy time series; and Pereira and Silveira
\cite{pereira2018unsupervised} use variational recurrent autoencoders with
attention. These avoid the labelling problem but differ markedly in
computational cost: deep autoencoder approaches are powerful yet heavy, whereas
the Isolation Forest is lightweight. Notably, Jamshidi \textit{et al.}
\cite{jamshidi2025lightweight} pair an autoencoder with an Isolation Forest
specifically for resource-constrained green IoT edge gateways, evidence that the
Isolation Forest is well matched to low-cost hardware. For an SME setting
characterised by both label scarcity and a tight computational budget, the
literature thus points clearly towards a lightweight, unsupervised detector.
 
\subsection{Low-Cost Embedded and Edge Computing}
The feasibility of the proposed approach rests on running useful analytics on
inexpensive hardware. The Industrial IoT and Industry~4.0 paradigm
\cite{munirathinam2020industry} frames this shift, and the emergence of TinyML
\cite{warden2019tinyml} has shown that machine-learning inference can be hosted
on low-power microcontrollers. Applied studies confirm this in practice: Ken
\textit{et al.} \cite{ken2024advancing} perform real-time ozone prediction on
cost-effective edge devices, and---most directly relevant---Baskar \textit{et al.}
\cite{baskar2024edge} deploy the Arduino Nano BLE Sense, the same platform
adopted in this project, for real-time industrial wastewater monitoring. This
validates the chosen hardware for environmental sensing at the edge. The
remaining gap is integrative: edge-ML studies typically address a single
classification or regression task, and rarely combine forecasting, anomaly
detection, and visualisation into one reproducible end-to-end system. Reliable
communication underpins all of this, as documented in surveys of smart grid
communication infrastructure \cite{fan2012smart,yan2012survey}.
 
\subsection{Synthesis, Research Gap, and Derived Parameters}
Taken together, the literature establishes four points. IoT-based monitoring is
mature but concentrated at residential and utility scales, with the more
complete systems relying on costly infrastructure. The LSTM is an effective
forecaster but must be validated against simple baselines and is rarely assessed
for deployability. Unsupervised, lightweight anomaly detection---the Isolation
Forest in particular---is the natural fit where labels are scarce and compute is
limited. Low-cost edge hardware, including the specific Arduino platform used
here, has been shown viable for real-world sensing. What remains absent is a
single, affordable, openly reproducible system that integrates all four for the
industrial SME. Ahmed \textit{et al.} \cite{ahmed2020machine} move toward such
integration for smart grid and renewable districts, but not at the low-cost,
SME-accessible point this project targets.
 
From this synthesis, a set of design parameters can be derived. These parameters
act as the bridge between the literature and the proposed solution architecture:
each is grounded in the evidence above and maps directly to a design decision,
as summarised in Table~\ref{tab:params}.
 
\begin{table*}[t]
\centering
\caption{Design parameters derived from the literature and their realisation in
the proposed system}
\label{tab:params}
\small
\renewcommand{\arraystretch}{1.3}
\begin{tabular}{|p{3.0cm}|p{9.4cm}|p{4.2cm}|}
\hline
\textbf{Parameter} & \textbf{Insight / evidence from the literature} &
\textbf{Realisation in this work}\\
\hline
Affordability &
Capable integrated systems rely on costly edge--cloud or proprietary
infrastructure \cite{bin2022real,silva2023iot}. &
Low-cost embedded platform (Arduino Nano BLE~33 Sense, $\sim$\pounds30).\\
\hline
Real-time operation &
Value comes from continuous feedback rather than periodic audits
\cite{hafeez2020efficient,bin2022real}. &
Streaming telemetry surfaced through a live dashboard.\\
\hline
Forecasting model &
LSTM captures multi-scale temporal patterns
\cite{hochreiter1997long,kong2017short}. &
Two-layer LSTM demand forecaster.\\
\hline
Baseline comparison &
ML does not automatically beat simple baselines
\cite{makridakis2018statistical}. &
Benchmarked against naive and moving-average baselines.\\
\hline
Anomaly paradigm &
Labelled faults are rarely available
\cite{himeur2021artificial,jesmeen2021unsupervised}. &
Unsupervised, label-free Isolation Forest.\\
\hline
Computational footprint &
Lightweight detectors suit constrained edge hardware
\cite{jamshidi2025lightweight}. &
Isolation Forest over heavy deep autoencoders.\\
\hline
Data-pipeline rigour &
Sound evaluation requires leakage-free protocols
\cite{makridakis2018statistical}. &
Documented preprocessing; chronological train/val/test split.\\
\hline
Reproducibility &
Few studies release a complete, replicable pipeline. &
Public code, processed data, and trained models.\\
\hline
Interpretability &
Black-box models hinder operator trust \cite{lundberg2017unified}. &
Provision for SHAP-based explanation (future work).\\
\hline
\end{tabular}
\end{table*}
 
\subsection{Proposed Solution Architecture}
Guided by these parameters, the proposed system is a four-layer architecture: an
IoT sensing layer built on the validated low-cost Arduino platform
\cite{baskar2024edge}; a data-pipeline layer providing documented, leakage-free
preprocessing; a machine-learning layer pairing an LSTM forecaster
\cite{hochreiter1997long,kong2017short} with an unsupervised Isolation Forest
detector \cite{liu2008isolation}; and a visualisation layer delivering a
real-time dashboard. Compared with the integrated edge--cloud system of Bin
Mofidul \textit{et al.} \cite{bin2022real}, the proposed design deliberately
trades some infrastructural sophistication for affordability and
reproducibility; and unlike the many studies that treat forecasting
\cite{wen2024deep} or anomaly detection \cite{pereira2018unsupervised} in
isolation, it integrates both within a single, openly released pipeline. The
detailed design and implementation of this architecture are presented in
Section~\ref{sec:method}.
\begin{figure}[t]
  \centering
  \includegraphics[width=\columnwidth]{images/Architecture_Professor_Format.png}
  \caption{Proposed four-layer system architecture.}
  \label{fig:arch}
\end{figure}
\section{Methodology and System Design}\label{sec:method}
\subsection{Research Approach}
The project follows a Design Science Research (DSR) methodology
\cite{hevner2004design,peffers2007design}, in which knowledge is generated
through the construction and rigorous evaluation of an artefact---here, the
integrated smart grid monitoring system. Because the requirements were
exploratory and the work was experimentation-heavy, an incremental,
experimental process was adopted, comprising four phases: data acquisition and
preparation, model development, evaluation against baselines, and integration of
the models into a monitoring dashboard. The physical sensing layer follows the
low-cost architecture set out in Section~\ref{sec:litreview}; the evaluation
reported here uses a public benchmark dataset as its data source, with live
hardware deployment reserved for future work.
 
\subsection{Dataset and Preprocessing}
The primary dataset is the UCI Individual Household Electric Power Consumption
dataset \cite{hebrail2012individual}, comprising 2{,}075{,}259 minute-level
measurements recorded between December 2006 and November 2010. It contains seven
electrical quantities: global active power, global reactive power, voltage,
global intensity, and three sub-metering channels. This dataset was selected
because it provides a long, continuous, real-world record whose daily, weekly,
and seasonal structure is well suited to sequence modelling. It is worth noting
that other public benchmarks target related but distinct problems; the Smart
Grid Stability dataset \cite{arzamasov2018towards}, for example, poses a binary
stability-classification task rather than continuous demand forecasting, and is
therefore outside the scope of the forecasting model developed here.
 
The raw data were prepared through a documented, version-controlled pipeline.
First, 25{,}979 missing values were imputed by forward-fill, the standard choice
for time-series data in which the most recent observation is the best estimate
of the next. Second, the series was resampled from one-minute to hourly
resolution, producing 34{,}589 clean records and reducing noise while preserving
the temporal patterns relevant to energy management. Third, extreme values were
screened using the interquartile-range method. Fourth, all features were
normalised to the range $[0,1]$ using min--max scaling, as required for stable
recurrent-network training. Fifth, for the forecasting task a sliding window of
24 hours was applied, yielding 34{,}565 supervised input--output sequences.
Finally, the data were partitioned chronologically into training, validation,
and test sets in a 70/15/15 ratio (24{,}195, 5{,}184, and 5{,}186 samples
respectively); the scaler was fitted on the training data only, so that no
information from future time steps could leak into the model input
\cite{makridakis2018statistical}. Each feature was scaled to the range $[0,1]$ using min--max normalisation,
 
\begin{equation}
x' = \frac{x - x_{\min}}{x_{\max} - x_{\min}}, \label{eq:minmax}
\end{equation}
 
where $x_{\min}$ and $x_{\max}$ are computed on the training partition only. The
forecasting task was then framed as a supervised sliding-window problem: given
the previous $w = 24$ hourly observations, the model predicts the next hour,
 
\begin{equation}
\hat{y}_{t} = f\!\left(x_{t-w}, x_{t-w+1}, \dots, x_{t-1}\right), \label{eq:window}
\end{equation}
 
where $f(\cdot)$ is the trained LSTM.
 
\subsection{Demand Forecasting Model}
Short-term demand forecasting was addressed with a Long Short-Term Memory (LSTM)
network \cite{hochreiter1997long}, whose gating mechanism is well matched to the
multi-scale temporal dependencies in energy data and which has an established
record in load forecasting \cite{kong2017short}. At each time step $t$, an LSTM unit updates its cell state $C_t$ and hidden
state $h_t$ from the current input $x_t$ and the previous hidden state
$h_{t-1}$ through three gates---forget ($f_t$), input ($i_t$), and output
($o_t$)---defined as
 
\begin{align}
f_t &= \sigma\!\left(W_f \cdot [h_{t-1}, x_t] + b_f\right), \label{eq:forget}\\
i_t &= \sigma\!\left(W_i \cdot [h_{t-1}, x_t] + b_i\right), \label{eq:input}\\
\tilde{C}_t &= \tanh\!\left(W_C \cdot [h_{t-1}, x_t] + b_C\right), \label{eq:cand}\\
C_t &= f_t \odot C_{t-1} + i_t \odot \tilde{C}_t, \label{eq:cell}\\
o_t &= \sigma\!\left(W_o \cdot [h_{t-1}, x_t] + b_o\right), \label{eq:output}\\
h_t &= o_t \odot \tanh(C_t), \label{eq:hidden}
\end{align}
 
where $W_{\{f,i,C,o\}}$ and $b_{\{f,i,C,o\}}$ are the learned weight matrices and
bias vectors, $\sigma(\cdot)$ is the logistic sigmoid, $\tanh(\cdot)$ is the
hyperbolic tangent, and $\odot$ denotes element-wise multiplication. The forget
gate \eqref{eq:forget} controls how much of the previous cell state is retained,
the input gate \eqref{eq:input} and candidate \eqref{eq:cand} determine what new
information is written, and the output gate \eqref{eq:output} produces the hidden
state passed to the next step. This gating is what allows the network to retain
information over the long daily and weekly cycles present in the energy series. The network comprises two
stacked LSTM layers of 64 and 32 units, each followed by a dropout layer with
rate 0.2, and a dense output neuron, giving 29{,}857 trainable parameters. It
was trained to minimise the mean-squared-error loss using the Adam optimiser
with a batch size of 32 for up to 50 epochs, with early stopping (patience of 10
epochs) to prevent overfitting; the implementation used TensorFlow/Keras
\cite{abadi2016tensorflow} with a fixed random seed for reproducibility. To place
the model's performance in context, two transparent baselines were implemented:
a naive persistence predictor, which forecasts the next hour as equal to the
current one, and a 24-hour moving average.
 
\subsection{Anomaly Detection Model}
Anomalous consumption was identified with an Isolation Forest
\cite{liu2008isolation}, chosen because it is unsupervised---and therefore
requires no labelled fault data, which is unavailable for this dataset---and
because it is computationally lightweight and thus suited to low-cost
deployment. The detector operated on all seven electrical features
simultaneously (a multivariate approach), after standardising each feature to
zero mean and unit variance so that no single quantity dominated by virtue of
its scale. The model was configured with 200 base estimators, a contamination
rate of 0.01, automatic subsampling, and a fixed random seed. Because the
dataset carries no ground-truth anomaly labels, a transparent statistical
baseline---a Z-score rule flagging any hour whose global active power lies more
than three standard deviations from the mean---was used as a reference point for comparison. The complete detection procedure is summarised in Algorithm~\ref{alg:iforest}. 
\begin{algorithm}[t]
\caption{Multivariate Anomaly Detection using Isolation Forest}
\label{alg:iforest}
\small
\begin{algorithmic}[1]

\Require Dataset $X\in\mathbb{R}^{n\times7}$, contamination rate $c$, number of trees $T$
\Ensure Binary labels $y\in\{0,1\}^{n}$

\State Standardize features:
      $Z=(X-\mu)/\sigma$

\State Train an Isolation Forest with $T$ trees using $Z$

\For{each sample $z_i\in Z$}
    \State Compute anomaly score $s_i$
\EndFor

\State Select threshold $\tau$ such that a fraction $c$ of samples satisfy $s_i<\tau$

\For{each score $s_i$}
    \If{$s_i<\tau$}
        \State $y_i \gets \textbf{anomaly}$
    \Else
        \State $y_i \gets \textbf{normal}$
    \EndIf
\EndFor

\Return $y$

\end{algorithmic}
\end{algorithm} 
 
\subsection{Evaluation Metrics}
Formally, these quantities are defined as the LSTM was compared against the naive and
moving-average baselines on identical data. For anomaly detection, in the
absence of ground-truth labels, precision, recall, and the F1-score were
computed \emph{relative to the Z-score reference}. These figures therefore
measure the degree of agreement between the two methods rather than correctness
against true anomalies; they are complemented by an analysis of the multivariate
anomalies that the Isolation Forest detects but the univariate baseline misses,
which characterises the added value of the multivariate approach. Forecasting accuracy was measured with the mean absolute error (MAE), the
root-mean-square error (RMSE), the mean absolute percentage error (MAPE), and
the coefficient of determination ($R^2$),
 
\begin{align}
\mathrm{MAE} &= \frac{1}{n}\sum_{t=1}^{n}\left|y_t - \hat{y}_t\right|, \label{eq:mae}\\
\mathrm{RMSE} &= \sqrt{\frac{1}{n}\sum_{t=1}^{n}\left(y_t - \hat{y}_t\right)^2}, \label{eq:rmse}\\
\mathrm{MAPE} &= \frac{100}{n}\sum_{t=1}^{n}\left|\frac{y_t - \hat{y}_t}{y_t}\right|, \label{eq:mape}\\
R^2 &= 1 - \frac{\sum_{t=1}^{n}(y_t-\hat{y}_t)^2}{\sum_{t=1}^{n}(y_t-\bar{y})^2}, \label{eq:r2}
\end{align}
 
where $y_t$ is the actual value, $\hat{y}_t$ the prediction, $\bar{y}$ the mean
of the actual values, and $n$ the number of test observations. For the anomaly
detector, agreement with the Z-score reference was quantified using precision,
recall, and their harmonic mean, the F1-score,
 
\begin{equation}
P = \frac{TP}{TP+FP}, \quad
R = \frac{TP}{TP+FN}, \quad
F_1 = \frac{2PR}{P+R}, \label{eq:prf}
\end{equation}
 
where $TP$, $FP$, and $FN$ are the counts of records flagged by both methods, by
the Isolation Forest only, and by the Z-score reference only, respectively. The
Z-score used by the baseline is
 
\begin{equation}
z_t = \frac{x_t - \mu}{\sigma}, \label{eq:zscore}
\end{equation}
 
with $\mu$ and $\sigma$ the mean and standard deviation of global active power;
an hour is flagged when $|z_t| > 3$.
 
\subsection{Implementation and Reproducibility}
The system was implemented in Python~3.11 using TensorFlow/Keras, scikit-learn,
pandas, and NumPy. Development and training were carried out on a consumer laptop
(AMD Ryzen~7 6800U, 16~GB RAM), with the option of GPU-accelerated training on a
cloud platform for larger workloads. The visualisation layer is a Streamlit
dashboard that loads the saved model outputs and presents demand forecasts
alongside anomaly alerts. To support full reproducibility, fixed random seeds
were used throughout, and all source code, the processed dataset, trained model
files, and results have been released in a public repository.

\section{Results and Discussion}\label{sec:results}
This section reports the performance of the two models on the held-out data and
then discusses their practical significance and limitations. Unless otherwise
stated, forecasting metrics are reported over a 500-hour window of the held-out
test set, consistent with the figures presented in the monitoring dashboard.
 
\subsection{Demand Forecasting Performance}
Table~\ref{tab:forecast} reports the accuracy of the LSTM forecaster alongside
the two baselines. The LSTM achieved the lowest error on the metrics that matter
most for this task, with a mean absolute error (MAE) of 0.3707~kW and a
root-mean-square error (RMSE) of 0.5421~kW. This corresponds to an average
prediction that is within roughly 370~W of the true hourly demand. Figure~\ref{fig:uciforecast} illustrates the forecast tracking the actual series. Relative to
the naive persistence baseline the LSTM reduced MAE by 8.6\%, and relative to the
24-hour moving average it reduced MAE by 35.7\%. The coefficient of
determination tells the same story: the LSTM explains 44.3\% of the variance in
hourly demand ($R^2 = 0.4432$), nearly double the naive baseline ($0.2361$) and
far above the moving average ($0.0542$).
 
\begin{table}[t]
\centering
\caption{Forecasting performance on the held-out test set (lower is better, except $R^2$)}
\label{tab:forecast}

\resizebox{\columnwidth}{!}{%
\begin{tabular}{|l|c|c|c|c|}
\hline
\textbf{Model} & \textbf{MAE (kW)} & \textbf{RMSE (kW)} & \textbf{MAPE (\%)} & \textbf{$R^2$} \\
\hline
Naive & 0.4054 & 0.6350 & 44.07 & 0.2361 \\
\hline
Moving Average & 0.5768 & 0.7066 & 87.89 & 0.0542 \\
\hline
\textbf{LSTM} & \textbf{0.3707} & \textbf{0.5421} & 45.77 & \textbf{0.4432} \\
\hline
\end{tabular}
}
\end{table}
\begin{figure}[t]
  \centering
  \includegraphics[width=\columnwidth]{images/04_actual_vs_predicted.png}
  \caption{LSTM forecast against actual demand on the UCI test set.}
  \label{fig:uciforecast}
\end{figure}
 
One result in Table~\ref{tab:forecast} warrants explicit comment. The LSTM's
mean absolute percentage error (MAPE) of 45.77\% is marginally higher than that
of the naive baseline (44.07\%), even though the LSTM is clearly superior on MAE,
RMSE, and $R^2$. This apparent contradiction arises from a well-known weakness of
MAPE on this type of data: household consumption falls close to zero during the
night, and dividing by these near-zero values inflates the percentage error
disproportionately, regardless of how small the absolute error is. MAPE is
therefore not a reliable discriminator here, and the absolute measures (MAE and
RMSE) and the variance-explained measure ($R^2$) give a more faithful picture of
performance. The moving average performs worst throughout, as its smoothing
lags the sharp peaks in demand that the LSTM is able to anticipate.
 
\subsection{Anomaly Detection Performance}
The Isolation Forest, operating on all seven electrical features, flagged 346 of
the 34{,}589 hourly records (1.00\%) as anomalous. Figure~\ref{fig:ucianom} shows their distribution across the series and in score space. To assess these detections in
the absence of ground-truth labels, they were compared against the univariate
Z-score baseline, which flagged 440 records. Table~\ref{tab:anomaly} summarises
the agreement between the two methods. The methods agreed on 267 records; the
Isolation Forest additionally flagged 79 records that the Z-score rule did not,
while the Z-score rule flagged 173 that the Isolation Forest did not.
 
\begin{table}[t]
\centering
\caption{Agreement between the Isolation Forest and the Z-score baseline}
\label{tab:anomaly}
\small
\renewcommand{\arraystretch}{1.3}
\begin{tabular}{|l|c|}
\hline
\textbf{Category} & \textbf{Records}\\
\hline
Flagged by Isolation Forest & 346\\
\hline
Flagged by Z-score baseline & 440\\
\hline
Flagged by both & 267\\
\hline
Isolation Forest only (multivariate) & 79\\
\hline
Z-score only (univariate) & 173\\
\hline
\end{tabular}
\end{table}
\begin{figure}[t]
  \centering
  \includegraphics[width=\columnwidth]{images/05_anomaly_detection.png}
  \caption{Detected anomalies across the UCI series (top) and the anomaly-score distribution (bottom).}
  \label{fig:ucianom}
\end{figure}
 
Treating the Z-score flags as a reference, the Isolation Forest achieved a
precision of 0.7717, a recall of 0.6068, and an F1-score of 0.6794. It is
important to interpret these figures correctly. Because no ground-truth anomaly
labels exist, they measure the \emph{agreement} between the two methods rather
than correctness against true anomalies; the Z-score rule is a transparent point
of reference, not an authoritative ground truth. Understood in this way, the
results indicate that the Isolation Forest reproduces the clear, high-magnitude
anomalies that a simple statistical rule would catch, while its principal value
lies in the 79 additional records it identifies. Inspection of these records
showed unusual \emph{combinations} of readings---for example, elevated intensity
and sub-metering activity without an extreme total-power value---that a
single-feature rule cannot detect by construction. This is the concrete benefit
of the multivariate approach anticipated in the literature review.
 
\subsection{Discussion}
Taken together, the results support the central claim of this work: that
capable energy analytics can be delivered with lightweight, low-cost methods.
The LSTM provides a genuine improvement over standard baselines on the metrics
appropriate to the task, and the Isolation Forest adds a complementary,
label-free capability for surfacing unusual consumption, including multivariate
patterns invisible to univariate rules. Both models are small enough to be
compatible with the affordable-deployment goal that motivates the project, and
the complete pipeline has been integrated into a working monitoring dashboard.
 
Several limitations should be acknowledged. The forecasting metrics are reported
over a bounded test window rather than the full multi-year series, and the
anomaly evaluation, lacking labelled faults, is necessarily an agreement study
rather than a measurement of true detection accuracy. The models were also
trained on a single residential dataset; although this dataset is a
well-established benchmark, industrial load profiles may differ, and validation
on a second, more industrial dataset is identified as the immediate next step to
strengthen external validity. Finally, both models are presently black boxes;
integrating post-hoc explanation would improve operator trust in a deployed
setting. These limitations are revisited in Section~\ref{sec:conclusion}.

\section{Conclusion and Future Work}\label{sec:conclusion}
This paper set out to demonstrate that intelligent energy management---real-time
monitoring, short-term demand forecasting, and anomaly detection---can be
delivered at a cost accessible to small and medium-sized enterprises, using
low-cost hardware and open-source software rather than proprietary
infrastructure. To that end, an end-to-end pipeline was designed,
implemented, and evaluated: a documented, leakage-free data-preparation
stage; an LSTM demand forecaster benchmarked against transparent baselines; a
multivariate Isolation Forest anomaly detector; and a real-time monitoring
dashboard, with the complete system released publicly for independent
replication.
 
The empirical results support the central claim. On the UCI household
benchmark, the LSTM achieved an MAE of 0.3707~kW, reducing error by 8.6\%
against a naive predictor and 35.7\% against a 24-hour moving average, while
the Isolation Forest surfaced 79 multivariate anomalies that a univariate
statistical rule cannot detect by construction. Cross-dataset validation on
the Tetouan City load confirmed that the unchanged pipeline transfers across
scales, with the forecaster's advantage widening on the smoother aggregated
series (62.0\% and 87.0\% improvements over the same baselines, nMAE of
2.36\%) and the anomaly detector shifting, as expected, towards
weather--demand combination anomalies. Together, the two evaluations indicate
that the proposed design meets its parameters of affordability, baseline-
validated accuracy, label-free anomaly detection, and reproducibility.
 
The work has limitations that define its future direction. First, the
evaluation, although spanning two public datasets at different scales, has not
yet included live industrial sensor data; the immediate next step is to deploy
the designed sensing layer on the Arduino Nano BLE~33 Sense platform and to
validate the pipeline on measurements collected in situ. Second, the anomaly
evaluation remains an agreement study in the absence of labelled faults;
deployment feedback from operators would allow detections to be confirmed and
the contamination parameter to be tuned against ground truth. Third, both
models are currently black boxes; integrating post-hoc explanation, such as
SHAP-based attribution, would let operators see \emph{why} an hour was
forecast high or flagged anomalous, improving trust in a deployed setting.
Finally, the forecaster could be extended from univariate to multivariate
input, exploiting the meteorological variables shown by the Tetouan analysis
to interact strongly with demand.
 
In sum, the contribution of this work is not a new algorithm but a
demonstration: that established, lightweight methods, engineered carefully
into a single reproducible system, are sufficient to bring smart-grid
capabilities within reach of the organisations that currently lack them.

\bibliographystyle{IEEEtran}
\bibliography{References}

\end{document}